\section{Instruction study: prompts, grading rubrics, and token accounting}
\label{app:cost-repro-procedures}
This appendix gives the inputs, the exact prompts, the execution settings, the grading rubrics, and the token accounting behind the instruction comparisons of Appendix~\ref{app:complementary-detail} and Table~\ref{tab:cost-repro}. The repository holds the saved inputs, answers, model grades, and usage records in \path{studies/instructions/}, and its \path{README.md} explains how to reproduce every table from them. The prompt quotations reproduce \path{studies/instructions/reproduction/prompts.json}, with only LaTeX escaping and dash formatting changed.

\subsection{Study design}
\label{app:cost-repro-reproduction}

\paragraph{Question and comparisons.}
The study asks how the instruction changes the effort needed to identify a false claim. The initial experiment applied four instructions to 100 proofs with Opus~5 at high effort in the Claude Code execution system; the reproduction ran the same 100 proofs with \texttt{gpt-5.6-sol} and Gemini~3.1~Pro, and a separate set of 60 statements without proofs with all three models, reusing the saved Opus~5 proof answers for its paired comparisons. Instructions are compared within each model, with the token accounting that its execution system provides.

\paragraph{The 100-proof input set.}
Each \emph{proofs100} record contains a problem statement and a numbered proof, stored in the JSON fields \texttt{record\_id}, \texttt{problem}, and \texttt{proof\_numbered}; a manifest fixes the input list and the input and reference files by SHA-256 hashes. These are the earlier instruction-study inputs, not the 100 benchmark inputs drawn from the human-approved core.

\paragraph{The 60-statement input set.}
The \emph{claims60} panel contains 60 statements assessed as false in a skeptical review, each with a reference counterexample. Executable checks support 57 of the references, 47 marked fully verified and 10 with exhaustive-search witnesses; the other three rest on the literature, recorded as two partial checks and one citation-only check in the \texttt{verification\_level} field. Each model receives the statement with its definitions, conventions, and an opaque identifier. The claims and witnesses are separate from BugBank records.

Some definitions hint at a refutation. Before the study ran, we fixed a text rule to identify such hints, the case-insensitive regular expression \texttt{/counterexample\textbar{}claim under test\textbar{}the claim (asserts\textbar{}quantifies)\textbar{}to be refuted\textbar{}refut/i} applied to the definitions field; it flags 13 of the 60 claims. Section~\ref{app:claims60-corrected-accounting} summarizes the comparisons on the remaining 47 no-hint claims and gives the complete sensitivity table.

\paragraph{Arms and prompts.}
Each cell sends one prompt. The proof template is

\begin{quote}\small\ttfamily\raggedright You are assessing a mathematical proof for correctness. Read \{input\_path\} --- it contains the problem statement and a line-numbered proof. Use ONLY that file: no web, no other files, no code execution. (Reading that file with your file-reading tool, or with a single shell command that only prints it, is allowed and is not code execution.)\par \{method\}\par Honesty rules: do not assume the proof is flawed; many proofs are fine. Do not pad. In your structured output set record\_id to exactly \{record\_id\_json\} and arm to exactly "\{arm\}".\end{quote}
with \texttt{\{method\}} one of the instructions below (the arm), \texttt{\{input\_path\}} the record's
input file, and \texttt{\{record\_id\_json\}} and \texttt{\{arm\}} the values the answer must echo.
\noindent\emph{\texttt{refute}}:
\begin{quote}\small\ttfamily\raggedright METHOD you must follow --- COUNTEREXAMPLE-FIRST: read the proof looking for its riskiest concrete claims, and try to construct ONE concrete counterexample (explicit numbers, explicit small structure, explicit trace) to some claim the proof actually relies on. The moment you have a counterexample you have checked carefully, report INVALID with that counterexample and STOP --- do not audit the rest of the proof. If after honest effort you cannot refute anything and believe the proof is correct, report VALID with counterexample "".\end{quote}
\noindent\emph{\texttt{verify\_full}}:
\begin{quote}\small\ttfamily\raggedright METHOD you must follow --- COMPLETE VERIFICATION: verify the proof step by step. For EVERY numbered step, decide whether it is justified given the problem and the preceding steps, and record a one-line justification or objection per step in assessment\_summary. You must judge every step even after finding a problem. Verdict INVALID if any step fails (name the first failing step in flawed\_step), VALID otherwise. counterexample may be "" --- your job is verification, not witness construction.\end{quote}
\noindent\emph{\texttt{verify\_earlystop}}:
\begin{quote}\small\ttfamily\raggedright METHOD you must follow --- VERIFICATION WITH EARLY STOP: verify the proof step by step in order. For each numbered step, decide whether it is justified given the problem and the preceding steps. You MAY stop at the first step you are certain fails; report INVALID with that step in flawed\_step. If you reach the end, report VALID. counterexample may be "" --- your job is verification, not witness construction.\end{quote}
\noindent\emph{\texttt{verify\_earlystop\_witness}}:
\begin{quote}\small\ttfamily\raggedright METHOD you must follow --- VERIFICATION WITH EARLY STOP, WITNESS REQUIRED: verify the proof step by step in order. For each numbered step, decide whether it is justified given the problem and the preceding steps. You MAY stop at the first step you are certain fails --- but a verdict of INVALID is only complete when you ALSO produce a CONCRETE counterexample (explicit numbers, explicit small structure, explicit trace) to that failing step, checked carefully, in the counterexample field. This arm exists to equalize deliverables: your INVALID must come with the same artifact a counterexample-first assessment produces. If you reach the end without a definitive failure, report VALID with counterexample "".\end{quote}
For the claims the template reads
\begin{quote}\small\ttfamily\raggedright You are assessing a mathematical statement for truth. Read \{input\_path\} --- it contains the statement and the definitions and conventions under which it is to be read. Use ONLY that file: no web, no other files, no code execution. (Reading that file with your file-reading tool, or with a single shell command that only prints it, is allowed and is not code execution.)\par \{method\}\par Honesty rules: do not assume the statement is false; many statements are true. Do not pad. In your structured output set record\_id to exactly \{record\_id\_json\} and arm to exactly "\{arm\}".\end{quote}
and the verification instruction becomes a complete proof attempt; the counterexample-first instruction is
unchanged in method:
\noindent\emph{\texttt{refute}}:
\begin{quote}\small\ttfamily\raggedright METHOD you must follow --- COUNTEREXAMPLE-FIRST: read the statement looking for its riskiest implication, and try to construct ONE concrete counterexample (explicit numbers, explicit small structure, explicit trace) that satisfies every hypothesis of the statement and violates its conclusion. The moment you have a counterexample you have checked carefully, report INVALID with that counterexample and STOP --- do not attempt a proof. If after honest effort you cannot refute it and believe the statement is true, report VALID with counterexample "".\end{quote}
\noindent\emph{\texttt{prove\_full}}:
\begin{quote}\small\ttfamily\raggedright METHOD you must follow --- COMPLETE PROOF ATTEMPT: attempt to prove the statement. Write a complete numbered step-by-step argument in assessment\_summary and, for EVERY numbered step, decide whether it is justified given the definitions and the preceding steps, recording a one-line justification or objection per step. You must judge every step even after finding a problem. Verdict VALID if every step is justified and the argument proves the statement; INVALID if the argument fails and you conclude the statement is false (name the first failing step in flawed\_step). counterexample may be "" --- your job is verification by proof, not witness construction.\end{quote}
\noindent\emph{\texttt{prove\_earlystop}}:
\begin{quote}\small\ttfamily\raggedright METHOD you must follow --- PROOF ATTEMPT WITH EARLY STOP: attempt to prove the statement step by step in order. For each numbered step, decide whether it is justified given the definitions and the preceding steps. You MAY stop at the first step you are certain cannot be justified; report INVALID with that step in flawed\_step. If your argument completes, report VALID. counterexample may be "" --- your job is verification by proof, not witness construction.\end{quote}
\noindent\emph{\texttt{prove\_earlystop\_witness}}:
\begin{quote}\small\ttfamily\raggedright METHOD you must follow --- PROOF ATTEMPT WITH EARLY STOP, WITNESS REQUIRED: attempt to prove the statement step by step in order. For each numbered step, decide whether it is justified given the definitions and the preceding steps. You MAY stop at the first step you are certain cannot be justified --- but a verdict of INVALID is only complete when you ALSO produce a CONCRETE counterexample (explicit numbers, explicit small structure, explicit trace) to the statement, checked carefully, in the counterexample field. This arm exists to equalize deliverables: your INVALID must come with the same artifact a counterexample-first assessment produces. If your argument completes without a definitive failure, report VALID with counterexample "".\end{quote}
Within each input set the four instructions share the surrounding prompt, the output schema, the grading rubric, and the grader; only the method block changes. The shared proof text has 338 characters before the method and 179 after it, and the statement text 367 and 186. Each condition receives one input file and no retrieved BugBank records.

The answer schema requires six string fields, \texttt{record\_id}, \texttt{arm}, \texttt{verdict}, \texttt{flawed\_step}, \texttt{counterexample}, and \texttt{assessment\_summary}, with the verdict one of \texttt{VALID}, \texttt{INVALID}, or \texttt{CANNOT\_DECIDE} and no extra fields. The initial study called the last field \texttt{reasoning\_summary}; the Opus~5 interface rejected that name during the reproduction, so all models used \texttt{assessment\_summary}.

\paragraph{Execution setup.}
Each model--input--instruction combination runs as one command-line session under Slurm (\url{https://slurm.schedmd.com/}) in a bubblewrap sandbox (\url{https://github.com/containers/bubblewrap}). The sandbox hides the research repository, scratch and temporary directories, and previous sessions. Within the session, it exposes the working directory, the answer schema, and one input file, with credentials supplied separately, while the rest of the host file system and the network remain visible. The prompt permits only reading the input, either with a file-reading tool or with a shell command that prints the file.

The configurations are:
\begin{itemize}
\item Opus~5: \texttt{claude -p}, high effort, only the \texttt{Read} tool,
with the CLI's JSON-schema output option.
\item GPT-5.6 Sol: \texttt{codex exec}, codex-cli 0.154.0, high reasoning
effort, approvals disabled, read-only sandbox, and the CLI's output schema.
\item Gemini~3.1~Pro: the High variant in Google Antigravity CLI 1.2.3,
high effort, JSON-schema output, and a 55-minute print timeout.
\end{itemize}

Codex and Antigravity retain tools that their interfaces do not let the driver remove, so a rule fixed before the runs flags any Codex shell command other than printing the input and any Antigravity tool step other than viewing the input or finishing. Such calls occur in one to three Codex cells per proof instruction, two to six Gemini cells per proof instruction, and up to three Gemini cells per statement instruction; their effect on the statement comparisons is reported below. Answers returned without reading the input are rejected and retried.

\paragraph{Retries and recovery.}
The driver runs six concurrent sessions for each Sol run and four for the Opus and Gemini runs and for each grading job, with a 3{,}600-second limit per session. Transport failures, timeouts, refusals, and invalid answer formats allow up to three attempts. A rate- or usage-limit signal (a structured signal from Claude, or matching CLI error messages from Codex and Antigravity) triggers a 20-minute pause without consuming an attempt. Six consecutive failures trigger a separate 20-minute pause, and three such pauses stop the driver. A cell that exhausts its attempts remains without an answer. Jobs relaunched after the recorded Codex interruption received an enlarged limit of six attempts.

The ledger records every attempt with its prompt hash. Slurm job identifiers in the lock files allow an interrupted run to resume while keeping its completed answers. The driver can also recover an answer from the newest usage record if recording was interrupted, a fallback that no reported run needed.

\paragraph{Token accounting.}
The main ratios use the output tokens of the accepted attempt: Claude's per-model \texttt{outputTokens}, Codex's \texttt{turn.completed} output, and Gemini's result usage, each including the engine's reported reasoning. Uncached and cached input come from the same receipts. Additional attempts, and attempts whose usage could not be recovered, are reported separately; when an interrupted shell failed to save a usage file, the available usage was recovered from the event log.

In the initial Opus study the agents also wrote their answers through a tool call, which added that text to output usage; the reproduction emits the structured answer once. This changes the absolute count by roughly one answer per cell in every condition and leaves the paired design intact.

\subsection{Token accounting for the 60-statement panel}
\label{app:claims60-corrected-accounting}

\begin{table}[t]
\centering\small
\begin{tabular}{lccc}
\toprule
Model & Supported S/O/W & Output S/W & Recorded total S/W \\
\midrule
Sol & 60/42/59 & \shortstack{0.593\\{\scriptsize [0.457, 0.791]}} & \shortstack{0.959\\{\scriptsize [0.928, 0.982]}} \\[3pt]
Gemini & 58/43/57 & \shortstack{0.739\\{\scriptsize [0.569, 0.889]}} & \shortstack{0.903\\{\scriptsize [0.755, 1.063]}} \\[3pt]
Opus & 60/60/60 & \shortstack{0.476\\{\scriptsize [0.240, 0.940]}} & \shortstack{0.849\\{\scriptsize [0.685, 0.983]}} \\[3pt]
\bottomrule
\end{tabular}

\caption{The separate 60-statement panel. S is search, O an ordinary early-stop proof attempt, and W a proof attempt that must supply a counterexample. Confirmed-objection counts use the same reference-assisted grader. Ratios compare accepted-attempt token sums. Recorded total tokens are uncached input, cached input, and output (with reported reasoning counted once within output), including reported Opus auxiliary usage. Brackets are pointwise paired 95\% intervals; output uses 2,000 resamples and recorded totals 20,000.}
\label{tab:claims60-evidence}
\end{table}


The accounting joins all 720 accepted usage records to their native receipts and completed ledger attempts, and every answer and grade to its record and instruction; all 120 input and reference files match the source manifest. Recomputing usage leaves the saved answers, grades, and output-token counts unchanged. The complete component table for all four instructions (\path{claims60_corrected_components.tex}), the full-panel comparisons (\path{RESULTS.json} and \path{README.md}), and the 3,018 source-file hashes are in \path{studies/instructions/corrected_claims60_accounting/}, with portable replay instructions in \path{studies/instructions/README.md}; the complete subset-sensitivity table is Table~\ref{tab:claims60-all-sensitivities} below.

\paragraph{Counting each component once.}
Sol's native input count already includes cached reads, whereas Gemini's \texttt{input\_tokens} holds only uncached input, so its cached reads must be added; an earlier Gemini summary had subtracted them. Adding the disjoint components gives search/witness totals of 2,480,272/2,745,952 and a ratio of $0.903$, correcting $0.872$, and changes the search/full-proof ratio from $0.978$ to $1.031$; output ratios are unchanged. Opus receipts report uncached input, cache creation, cache reads, and output separately. Reasoning is counted once, within output, for every model.


Opus~5 reports usage for the target model and for auxiliary models. Its target-model search/witness totals are 751,620/895,394, a ratio of $0.839$; including the reported auxiliary usage gives 824,253/970,902 and $0.849$, the all-reported total that Table~\ref{tab:claims60-evidence} uses, while the saved accounting also keeps the target-model view. Sol and Gemini supply native aggregate totals without an auxiliary-model breakdown.


\paragraph{Do hints or additional executions change the pattern?}
Two observational subsets check the pattern: the no-hint subset removes the 13 claims flagged by the pre-run text rule, and the no-additional-execution subset keeps a task pair only when both of its runs recorded no extra execution. Neither subset randomizes tool permissions, and Opus has no execution-count field. For Gemini, excluding the pairs with additional executions changes the output and the all-reported total-token ratios as follows. Against a full proof attempt (59 pairs), the output ratio moves from $0.858$ to $0.844$ and the total-token ratio from $1.031$ to $0.968$. Against ordinary early stopping (59 pairs), they move from $1.275$ to $1.352$ and from $1.210$ to $1.046$. Against early stopping with a counterexample required (57 pairs), they move from $0.739$ to $0.654$ and from $0.903$ to $0.923$. Sol has no such pairs.

On the 47 no-hint claims, the search/witness output ratios are $0.535$, $0.689$, and $0.396$ for Sol, Gemini, and Opus, and the corresponding all-reported total-token ratios are $0.951$, $0.942$, and $0.809$: lower output-token counts can coexist with similar recorded totals. Table~\ref{tab:claims60-all-sensitivities} gives all three comparisons across all subsets, including the extra output used by search relative to ordinary early stopping.

\begingroup\scriptsize\setlength{\tabcolsep}{3pt}
\begin{longtable}{lllrrrr}
\caption{Claims60 instruction and subset sensitivity. S denotes search; F full proof attempt, O ordinary early stop, and W early stop requiring a witness. Counts use preserved grades; token ratios are S/comparator. Tool-free means no recorded additional execution in either paired cell. Exclusions are descriptive. Opus tool use is unrecorded, so no tool-free row is inferred. Output is the target model; recorded total tokens include every model reported in the accepted receipt.}\label{tab:claims60-all-sensitivities}\\
\toprule Model & Arm & Subset & $n$ & Catches S/C & Output & Recorded total \\
\midrule\endfirsthead
\toprule Model & Arm & Subset & $n$ & Catches S/C & Output & Recorded total \\
\midrule\endhead
Sol & F & All & 60 & 60/59 & 0.544 & 0.954 \\
Sol & F & No hint & 47 & 47/47 & 0.507 & 0.949 \\
Sol & F & Tool-free & 60 & 60/59 & 0.544 & 0.954 \\
Sol & F & Both exclusions & 47 & 47/47 & 0.507 & 0.949 \\
Sol & O & All & 60 & 60/42 & 1.347 & 1.015 \\
Sol & O & No hint & 47 & 47/36 & 1.260 & 1.011 \\
Sol & O & Tool-free & 60 & 60/42 & 1.347 & 1.015 \\
Sol & O & Both exclusions & 47 & 47/36 & 1.260 & 1.011 \\
Sol & W & All & 60 & 60/59 & 0.593 & 0.959 \\
Sol & W & No hint & 47 & 47/46 & 0.535 & 0.951 \\
Sol & W & Tool-free & 60 & 60/59 & 0.593 & 0.959 \\
Sol & W & Both exclusions & 47 & 47/46 & 0.535 & 0.951 \\
Gemini & F & All & 60 & 58/42 & 0.858 & 1.031 \\
Gemini & F & No hint & 47 & 46/37 & 0.768 & 1.028 \\
Gemini & F & Tool-free & 59 & 57/42 & 0.844 & 0.968 \\
Gemini & F & Both exclusions & 46 & 45/37 & 0.742 & 0.947 \\
Gemini & O & All & 60 & 58/43 & 1.275 & 1.210 \\
Gemini & O & No hint & 47 & 46/36 & 1.130 & 1.235 \\
Gemini & O & Tool-free & 59 & 57/43 & 1.352 & 1.046 \\
Gemini & O & Both exclusions & 46 & 45/36 & 1.198 & 1.027 \\
Gemini & W & All & 60 & 58/57 & 0.739 & 0.903 \\
Gemini & W & No hint & 47 & 46/45 & 0.689 & 0.942 \\
Gemini & W & Tool-free & 57 & 56/55 & 0.654 & 0.923 \\
Gemini & W & Both exclusions & 45 & 45/44 & 0.610 & 0.915 \\
Opus & F & All & 60 & 60/60 & 0.888 & 0.983 \\
Opus & F & No hint & 47 & 47/47 & 0.845 & 0.977 \\
Opus & O & All & 60 & 60/60 & 1.355 & 1.043 \\
Opus & O & No hint & 47 & 47/47 & 1.275 & 1.034 \\
Opus & W & All & 60 & 60/60 & 0.476 & 0.849 \\
Opus & W & No hint & 47 & 47/47 & 0.396 & 0.809 \\
\bottomrule
\end{longtable}

\endgroup



\paragraph{Additional attempts.}
The solver ledgers contain 730 attempts: Sol has no rejected attempt, Gemini has three with known usage totaling 29,135 tokens, and Opus has six failed attempts with unknown usage and one rate-limited attempt with known target-model usage of 283,653 tokens. Unknown usage is recorded as unknown, including when a receipt lists zero API messages. The accepted-attempt ratios therefore describe completed solver answers; grading, later checking, and unrecorded interrupted attempts need separate accounting before an end-to-end cost can be estimated.

\paragraph{Grading.}
The grader, \texttt{gpt-6-astra} through Codex CLI at maximum effort, grades the new answers and regrades the 400 original Opus proof answers, receiving the assessment and the held-out reference inline in the same read-only sandbox with its shell tool available. The recorded September~15 pause contains 1,700 completed grading cells, six with shell commands including three attempted arXiv downloads; the completed dataset contains 1,920 cells, seven with shell commands. The extra 220 completions and the additional command-bearing cell are documented in \path{studies/instructions/reproduction/grader_cohort.json}, with usage and recorded commands in its adjacent \path{grade_usage/} directory; the 1,700-cell pause cohort was reconstructed from timestamps.

The two saved rubrics follow; their judgments are weighed against the source and witness checks in Section~\ref{app:reference-reconciliation}. The proof rubric is:

\begin{quote}\small\ttfamily\raggedright You are grading one proof assessment against held-out ground truth.\par Read these two files:\\ 1. \{assessment\_path\} - the assessment (verdict, claimed flawed step, counterexample). If this file does not exist, use grade MISSED with justification "assessment absent (quota interruption)".\\ 2. \{truth\_path\} - the ground truth: the confirmed flaw and its independently verified witness.\par Every proof in this sample IS flawed (pipeline-confirmed). Grade the assessment:\\ - CAUGHT\_MATCHING: it reported INVALID and its flawed step/counterexample targets the SAME defect as the ground truth (rephrasing and a different-but-valid witness for the same claim both count as matching).\\ - CAUGHT\_DIFFERENT: it reported INVALID with a DIFFERENT flaw, and after checking its support concretely yourself, that flaw is genuine.\\ - MISSED: it reported VALID or CANNOT\_DECIDE.\\ - FALSE\_CATCH: it reported INVALID but the stated flaw does not hold up under your concrete check (wrong math, misread claim, invalid witness).\\ Be strict about FALSE\_CATCH vs CAUGHT\_DIFFERENT: verify, do not take the assessment's word.\par Set record\_id exactly \{record\_id\_json\}, arm "\{arm\}".\end{quote}
The statement rubric adds two outcomes, a correct verdict without a checkable refutation and an object the grader cannot verify; both input sets share this grade list, and in Sol's early-stop proof condition one answer received each added grade, counted separately from confirmed and unsupported objections. The exact statement rubric is:
\begin{quote}\small\ttfamily\raggedright You are grading one assessment of a mathematical statement against held-out ground truth.\par Read these two files:\\ 1. \{assessment\_path\} - the assessment (verdict VALID = the statement is true, INVALID = false, CANNOT\_DECIDE; a claimed failing step; a counterexample).\\ 2. \{truth\_path\} - the ground truth: the statement is FALSE and this file gives an independently verified counterexample (the reference witness) and, where recorded, why the statement fails.\par Every statement in this sample IS false (reviewer-confirmed, witness verified). Grade the assessment:\\ - CAUGHT\_MATCHING: it reported INVALID and its counterexample (or, if it gave no explicit witness, its stated reason) is a VALID refutation of the same kind as the reference (the same object or family, or the same failing mechanism; a different-but-valid witness of the same kind counts as matching).\\ - CAUGHT\_DIFFERENT: it reported INVALID with a DIFFERENT valid refutation (a genuinely different counterexample object or a different correct reason), and after checking its support concretely yourself, that refutation is genuine.\\ - MISSED: it reported VALID or CANNOT\_DECIDE.\\ - FALSE\_CATCH: it reported INVALID but the stated counterexample or reason does not hold up under your concrete check (wrong math, an object that violates a hypothesis or satisfies the conclusion, a misread statement, or no checkable support at all).\\ - CAUGHT\_UNSUPPORTED: it reported INVALID (the correct verdict) but gave neither a counterexample nor a self-contained argument that the statement is false. A 'stated reason' counts as a refutation only if it is itself a complete argument (for example a parity or counting argument, or a named object together with why it meets every hypothesis and fails the conclusion). 'My proof attempt fails at step k', 'the hypotheses do not force the conclusion', or naming a gap in a would-be proof is NOT a refutation.\\ - UNVERIFIED: it reported INVALID with a named or large object (a famous graph, a structure with dozens of vertices, a construction you cannot check by hand from the assessment's own text) that is NOT the reference object and that you cannot verify concretely. Never guess: use UNVERIFIED rather than CAUGHT\_DIFFERENT or FALSE\_CATCH in that case. If the object IS the reference object (or an explicitly stated variant the truth file already covers), grade CAUGHT\_MATCHING.\\ Be strict about FALSE\_CATCH vs CAUGHT\_*: verify concretely, do not take the assessment's word; a counterexample must satisfy every hypothesis in the definitions and violate the conclusion. FALSE\_CATCH is only for INVALID whose supplied witness or argument is wrong under your concrete check. Clauses in the definitions that describe what a counterexample must include are guidance to the solver: grade on whether the object refutes the statement under the stated conventions, performing any omitted verification (trace, tie-break check, non-isomorphism) yourself; omission of demanded verification text alone never makes a FALSE\_CATCH.\par Set record\_id exactly \{record\_id\_json\}, arm "\{arm\}".\end{quote}
